\section{Problem Formulation}\label{sec:problem}
Let $c$ specify a scene and instruction, $h_t$ the observation and controller history, and $Y(\tau)\in\{0,1\}$ the complete-task outcome of trajectory $\tau$. The evaluation objective is
\begin{equation}
J(\theta)=\mathbb{E}_{c,\tau\sim P_\theta(\cdot\mid c)}[Y(\tau)].
\end{equation}
Training may use simulator-derived intermediate predicates, but final evaluation uses the declared task verifier. We distinguish the exploration policy used for training from the deployment policy used for evaluation.

The CDPR action $a_t\in[-1,1]^5$ controls Cartesian translation, yaw, and the gripper. The policy receives camera observations, proprioception, and an instruction; simulator object geometry is used to calculate rewards and evaluation events. Episodes begin with an empty gripper and require the entire manipulation sequence. Strict success requires a physical grasp and held lift, intentional release, and native placement, without latched carry-slip or wrong-placement events. Native placement alone is therefore an insufficient endpoint. The exact event predicates, temporal latches, and abort handling are part of the evaluation contract.

LIBERO success is the environment's binary goal-verification outcome~\cite{libero}. Its phase tracker uses simulator state during training. This endpoint and the CDPR strict endpoint are reported separately; their percentages cannot be pooled into a common success rate. Similarly, repeated development scenes measure performance on a declared distribution rather than unseen-task generalization.
